% !TeX program = lualatex
% !TeX encoding = utf8
% !TeX spellcheck = uk_UA
% !TeX root = ../main.tex

%% --------------------------------------------------------
%% Chapter file: Fusion.tex
%% --------------------------------------------------------

\graphicspath{{\currfilebase/Pictures}}

\chapter{Керований термоядерний синтез}\label{sec:fusion}\hypertarget{fusion}{}
%% --------------------------------------------------------

Основний практичний інтерес фізики плазми пов'язаний із~проблемою \emphx{керованого
термоядерного синтезу} (КТС). Для достатньо інтенсивних реакцій синтезу речовину треба
нагріти до $\sim10$--$20$~кілоелектронвольт (десятки мільйонів кельвінів), а~за таких температур вона є плазмою.



\section{Реакції та паливо}



Найлегше здійснюється реакція дейтерію з~тритієм
% ~~~~~~~~~~~~~~~~~~~~~~~~~~~~~~~~~~~~~~~~~~~~~~~~~~~~~~~~
\begin{equation*}
	\mathrm D+\mathrm T\to{}^4\mathrm{He}\,(3{,}5~\text{МеВ})+n\,(14{,}1~\text{МеВ}),
\end{equation*}
% ~~~~~~~~~~~~~~~~~~~~~~~~~~~~~~~~~~~~~~~~~~~~~~~~~~~~~~~~
з~повним виходом $E_f=17{,}6$~МеВ. Швидкість реакцій характеризує величина $\langle\sigma
v\rangle$ (усереднена за максвеллівським розподілом): для D--T вона дорівнює приблизно
$1{,}4\times10^{-17}$, $1{,}1\times10^{-16}$, $4{,}3\times10^{-16}$,
$6{,}7\times10^{-16}~\mathrm{см}^3/\mathrm{с}$ при $T=5$, $10$, $20$,
$30$~кеВ. Дейтерій міститься в~морській воді (приблизно один атом на 6400 атомів водню;
загальна маса в~океанах $\sim4\times10^{13}$~т, що за порядком величини відповідає
$10^{20}$~кВт$\cdot$рік, тоді як споживана людством потужність становить
$\sim10^{10}$~кВт). Тритій радіоактивний (період піврозпаду $\approx12$~років) і~в~природі
практично відсутній, тому його треба виробляти в~самому реакторі: нейтрони реакції
опромінюють літієвий шар (\enquote{бланкет}), $n+{}^6\mathrm{Li}\to\mathrm
T+{}^4\mathrm{He}$. Реакція D--D потребує значно вищих температур ($\sim50$~кеВ).



\section{Чому потрібна саме така температура: бар'єр і~тунелювання}\label{sec:gamow}



Два ядра мають злитися на відстані $r\sim4$~фм (радіус дії ядерних сил), але вони відштовхуються
кулонівською силою. Висота кулонівського бар'єра для D--T
$U_C=Z_1Z_2e^2/r\approx1{,}44~\text{МеВ$\cdot$фм}/4~\text{фм}\approx0{,}4$~МеВ, тобто в~$\sim40$ разів
більше за $10$~кеВ. Класично злиття можливе лише для частинок із~далекого максвеллівського
«хвоста» $e^{-E/T}\sim e^{-40}$~--- замало. Реакції ідуть завдяки
\emph{квантовому тунелюванню}.

\begin{Theory}[Нагадаємо з~квантової механіки]
У~квазікласичному (ВКБ) наближенні ймовірність проходження крізь бар'єр дорівнює
$\exp\bigl(-\tfrac2\hbar\int_{r_1}^{r_2}|p|\,dr\bigr)$, де $|p|=\sqrt{2m_r(U(r)-E)}$ під бар'єром.
Для кулонівського бар'єра $U=Z_1Z_2e^2/r$ інтеграл від точки зупинки $r_2=Z_1Z_2e^2/E$ до ядерного
радіуса (при $E\ll U_C$ нижня межа несуттєва) обчислюється точно й~дає $\sqrt{E_G/E}$.
\end{Theory}

Ймовірність проходження бар'єра (формула Гамова)
% ~~~~~~~~~~~~~~~~~~~~~~~~~~~~~~~~~~~~~~~~~~~~~~~~~~~~~~~~
\begin{equation}\label{eq:gamow}
	P_{\rm tun}\approx\exp\!\left(-\sqrt{\frac{E_G}{E}}\right),\qquad
	E_G=2m_rc^2\,(\pi\alpha Z_1Z_2)^2,
\end{equation}
% ~~~~~~~~~~~~~~~~~~~~~~~~~~~~~~~~~~~~~~~~~~~~~~~~~~~~~~~~
де $m_r$~--- наведена маса, $\alpha=e^2/\hbar c$. Для D--T $m_rc^2=1{,}12$~ГеВ, $E_G\approx1{,}18$~МеВ.
Перерізи реакції тому швидко зростають з~енергією, $\sigma(E)=S(E)\,E^{-1}e^{-\sqrt{E_G/E}}$,
де $S(E)$ повільно змінюється (а~для D--T має резонанс при $\sim64$~кеВ, що додатково підсилює реакцію).

Швидкість реакції $\langle\sigma v\rangle$ визначається добутком двох протилежних факторів
(рис.~\ref{fig:gamow}): максвеллівського хвоста $e^{-E/T}$ (спадає) і~імовірності
тунелювання (зростає). Добуток має максимум~--- \emph{пік Гамова}~--- при
% ~~~~~~~~~~~~~~~~~~~~~~~~~~~~~~~~~~~~~~~~~~~~~~~~~~~~~~~~
\begin{equation}\label{eq:gamow-peak}
	E_0=\left(\frac{E_G\,(T)^2}{4}\right)^{1/3}\approx31~\text{кеВ при}\ T=10~\text{кеВ},
\end{equation}
% ~~~~~~~~~~~~~~~~~~~~~~~~~~~~~~~~~~~~~~~~~~~~~~~~~~~~~~~~~~~~~~
Максимум знаходимо з~умови $\dfrac{d}{dE}\Bigl(-\dfrac E{T}-\sqrt{\dfrac{E_G}E}\Bigr)=0$,
тобто $\dfrac1{T}=\dfrac12\sqrt{E_G}\,E^{-3/2}$, звідки $E_0=(E_GT^2/4)^{1/3}$.
У~максимумі $E_0/T=(E_G/4T)^{1/3}$, а~$\sqrt{E_G/E_0}=2(E_G/4T)^{1/3}$, тож
показник дорівнює $-3(E_G/4T)^{1/3}$. Фізично: реагують не типові частинки, а~хвіст
Максвелла~--- компроміс між бар'єром і~числом частинок. Отже,
реакції відбуваються переважно за участі частинок з~енергіями втричі вищими за $T$, а~
\begin{equation*}
	\langle\sigma v\rangle\propto T^{-2/3}\exp\!\left[-3\Bigl(\frac{E_G}{4T}\Bigr)^{1/3}\right]
	\ \Rightarrow\ \langle\sigma v\rangle\propto T^{2\text{--}3}\ \text{при}\ T\sim10~\text{кеВ}.
\end{equation*}
Саме тому швидкість реакції \emph{різко} (степенево, як $T^2$--$T^3$) зростає з~температурою в~області
кількох кеВ і~сильно спадає при $T\ll10$~кеВ. Крім того, $E_G\propto(Z_1Z_2)^2m_r$, тому для ядер із~більшими зарядами бар'єр вищий і~потрібні ще вищі температури. Для D--D $E_G\approx1{,}0$~МеВ, але в~цієї реакції немає резонансу, а~її $S$-фактор на два порядки менший, тому вона потребує $T\sim50$~кеВ.

\begin{figure}[h!]\centering
	\localinput{gamow.tikz}
	\caption{Пік Гамова для реакції D--T: добуток спадного максвеллівського хвоста та зростаючої
	ймовірності тунелювання дає максимум біля $E_0\approx31$~кеВ при $T=10$~кеВ}
	\label{fig:gamow}
\end{figure}

Залежність $\langle\sigma v\rangle(T)$ для основних реакцій показано на рис.~\ref{fig:sigmav}
(параметризація Боша~--- Гейла). Крива D--T лежить на кілька порядків вище за інші в~області
$T\lesssim100$~кеВ і~має широкий максимум поблизу $60$--$100$~кеВ; саме тому для першого реактора
обрано D--T і~температури $10$--$20$~кеВ: там $\langle\sigma v\rangle$ уже достатньо велике, а~втрати ще помірні.

\begin{figure}[ht]\centering
	\localinput{sigmav.tikz}
	\caption{Швидкість реакції $\langle\sigma v\rangle$ для D--T, D--$^3$He та D--D залежно від температури
	(логарифмічні осі, параметризація Боша~--- Гейла). Пунктиром позначено робочий інтервал $10$--$20$~кеВ}
	\label{fig:sigmav}
\end{figure}


\section{Баланс потужностей і~критерій Лоусона}



Нехай $n_D=n_T=n/2$, $n_e=n$, $T_e=T_i=T$. Потужність синтезу в~одиниці об'єму
$P_f=n_Dn_T\langle\sigma v\rangle E_f=\tfrac14n^2\langle\sigma v\rangle E_f$. Нейтрони
залишають плазму, а~заряджені $\alpha$-частинки ($E_\alpha=3{,}5$~МеВ) передають їй свою
енергію: $P_\alpha=\tfrac14n^2\langle\sigma v\rangle E_\alpha$. Тепловий запас плазми
$W=3nT$ (по $\tfrac32T$ на електрон і~на іон), а~втрати
характеризуються часом утримання енергії $\tau_E$: $P_L=W/\tau_E$. Умова
\emphx{запалювання} (самопідтримна реакція) $P_\alpha\geq P_L$ дає
% ~~~~~~~~~~~~~~~~~~~~~~~~~~~~~~~~~~~~~~~~~~~~~~~~~~~~~~~~
\begin{equation}\label{eq:lawson}
	\tcbhighmath{n\tau_E\geq\frac{12\,T}{\langle\sigma v\rangle E_\alpha}}.
\end{equation}
% ~~~~~~~~~~~~~~~~~~~~~~~~~~~~~~~~~~~~~~~~~~~~~~~~~~~~~~~~
З~наведених значень $\langle\sigma v\rangle$: $n\tau_E\gtrsim3\times10^{14}$ при $10$~кеВ
і~мінімум $\approx1{,}5\times10^{14}~\mathrm{см}^{-3}\mathrm{с}$ при $20$--$30$~кеВ.
Помноживши на $T$, одержуємо потрійний добуток
$nT\tau_E\gtrsim3\times10^{15}~\text{кеВ}\cdot\mathrm{с}\cdot\mathrm{см}^{-3}$ (мінімум
поблизу $15$~кеВ), який і~прийнято вважати мірою досягнення. Звідки мінімум за $T$: $n\tau\propto T/\langle\sigma v\rangle$, тому мінімум відповідає
максимуму $\langle\sigma v\rangle/T$, а~не $\langle\sigma v\rangle$; тому мінімум $n\tau$ лежить при $\sim25$~кеВ, а~мінімум $nT\tau\propto T^2/\langle\sigma v\rangle$ при $\sim14$~кеВ.
Студенти плутають ці оптимуми. Це \emphx{критерій Лоусона}
в~сучасному формулюванні (оригінал 1957~р. стосувався беззбитковості при заданому ККД,
тому числові множники різняться). Для D--D, за оцінками, потрібне на два порядки більше:
$n\tau\sim10^{16}~\mathrm{см}^{-3}\mathrm{с}$.



\section{Втрати}



У~розрахунку $\tau_E$ не було враховано випромінювання. Основні радіаційні втрати~---
\emphx{гальмівне} (при кулонівських зіткненнях електронів з~іонами) і~\emphx{циклотронне}
(магнітогальмівне, від руху електронів по ларморівських колах). Потужність гальмівного
випромінювання $P_{br}\approx1{,}7\times10^{-32}\,n_e\sqrt{T_e}\sum
Z^2n_i~\mathrm{Вт/см}^3$ ($T_e$ в~еВ, $n$ в~$\mathrm{см}^{-3}$) зростає як $Z^2$, тому
домішки важких елементів небезпечні. Для $n=10^{14}~\mathrm{см}^{-3}$, $10$~кеВ маємо
$P_f\approx0{,}80$, $P_\alpha\approx0{,}16$, $P_{br}\approx0{,}017~\mathrm{Вт/см}^3$,
тобто гальмівне випромінювання становить близько 10\,\% від $\alpha$-нагрівання.
Враховуючи його, нижня межа температури запалювання D--T становить кілька кеВ, але
практично працюють при $10$--$20$~кеВ. Циклотронне випромінювання $\propto B^2nT$, проте
воно в~основному реабсорбується плазмою та відбивається стінками.



\section{Дві стратегії: магнітне й~інерційне утримання}



Умову~\eqref{eq:lawson} можна виконати за добутком $n\tau$ двома способами.
У~\emphx{магнітному утриманні} (розд.~\ref{sec:magnetic-confinement})
$n\sim10^{14}~\mathrm{см}^{-3}$ і~$\tau_E\sim$ секунди. У~\emphx{інерційному} стисканні
мішені лазерним або іншим випромінюванням густина досягає
$10^{25}$--$10^{26}~\mathrm{см}^{-3}$, а~реакція встигає відбутися за $\sim10^{-10}$~с,
поки речовина не розлетілася.



\section{Що означає цей критерій на практиці}



Критерій Лоусона не є самодостатньою умовою працездатності реактора: він не враховує всі
види випромінювання, нестійкості, транспорт і~технологічні втрати. Тому його використовують
як базовий орієнтир, а~не як гарантію. Для порівняння установок зручніше користуватися саме
потрійним добутком $nT\tau_E$ й~коефіцієнтом підсилення плазми $Q$ (див.~нижче).



\section{Від реакції до розетки}



Критерій Лоусона відповідає на питання, чи \emph{самопідтримується} реакція. Але мета~---
електрика. Простежмо, куди дівається енергія.



\subsection*{Куди дівається 17,6~МеВ}



Продуктів реакції два: $\alpha$-частинка та~нейтрон. Початковий сумарний імпульс
(у~системі центра мас) нульовий, тому за~збереженням імпульсу $p_\alpha=p_n$. Кінетична
енергія $p^2/2m$ обернено пропорційна масі, отже
% ~~~~~~~~~~~~~~~~~~~~~~~~~~~~~~~~~~~~~~~~~~~~~~~~~~~~~~~~
\begin{equation*}
	\frac{E_\alpha}{E_n}=\frac{m_n}{m_\alpha}\approx\frac14
	\quad\Rightarrow\quad
	E_\alpha=\frac{17{,}6}{5}\approx3{,}5~\text{МеВ},\qquad E_n\approx14{,}1~\text{МеВ}.
\end{equation*}
% ~~~~~~~~~~~~~~~~~~~~~~~~~~~~~~~~~~~~~~~~~~~~~~~~~~~~~~~~
Тобто $20\%$ енергії несе $\alpha$-частинка, $80\%$~--- нейтрон. Їхні долі різні, бо різні
маси: за однакового за модулем імпульсу легша частинка має більшу кінетичну енергію.

\emph{$\alpha$-частинка} заряджена, тому прив'язана до силових ліній. Її швидкість
$v_\alpha=\sqrt{2E_\alpha/m_\alpha}\approx1{,}3\times10^{9}$~см/с, і~в~полі $5$~Тл
ларморівський радіус $r_{L\alpha}=m_\alpha cv_\alpha/(2eB)\approx5$~см, тобто приблизно на
порядок менший за характерний розмір плазми близько метра. Це означає, що
$\alpha$-частинка є замагніченою, хоча сама по собі мала орбіта ще не гарантує утримання
в~усій конфігурації. У~добре замкненому магнітному полі вона може залишатися в~плазмі
достатньо довго, щоб віддати їй значну частину енергії: спершу електронам (її швидкість
менша за~теплову швидкість електронів, але більша за іонну), потім іонам. Це
і~є~$P_\alpha$ в~умові запалювання~\eqref{eq:lawson}.

Гальмування $\alpha$-частинок характеризує критична енергія
$$E_{\rm crit}\approx14{,}8\,T_eA_\alpha\Bigl(\sum_j\frac{n_jZ_j^2}{n_eA_j}\Bigr)^{2/3}\approx33\,T_e\approx330~\text{кеВ}\quad(T_e=10~\text{кеВ}).$$
Оскільки $E_\alpha=3{,}5$~МеВ $\gg E_{\rm crit}$, $\alpha$-частинка спочатку віддає енергію переважно електронам: її швидкість
$1{,}3\times10^{9}$~см/с менша за $v_{Te}\approx4\times10^{9}$~см/с, але набагато більша за $v_{Ti}$, тож застосовне динамічне тертя з~розд.~\ref{sec:plasma-collisions}~---
приклад того, що теорія зіткнень працює для реактора.

\emph{Нейтрон} не має заряду, магнітне поле на~нього не діє. Його вільний пробіг у~плазмі
$\sim1/(n\sigma)\sim10^{10}$~см (порядку~$10^{5}$~км) нескінченно більший за~установку,
тож нейтрон безперешкодно вилітає й~потрапляє у~\emph{бланкет}~--- шар завтовшки близько
метра навколо плазми.

\begin{figure}[h!]\centering
	\localinput{energy-flow.tikz}
	\caption{Потоки енергії в~термоядерній електростанції}
	\label{fig:energy-flow}
\end{figure}



\subsection*{Бланкет: тепло і~тритій}



У~бланкеті нейтрон гальмується в~зіткненнях, і~його енергія переходить у~тепло. Крім того,
він виконує другу, не менш важливу роботу: наробляє тритій для палива:
% ~~~~~~~~~~~~~~~~~~~~~~~~~~~~~~~~~~~~~~~~~~~~~~~~~~~~~~~~
\begin{equation*}
	n+{}^6\mathrm{Li}\to\mathrm T+{}^4\mathrm{He}+4{,}8~\text{МеВ},\qquad
	n+{}^7\mathrm{Li}\to\mathrm T+{}^4\mathrm{He}+n'-2{,}5~\text{МеВ}.
\end{equation*}
% ~~~~~~~~~~~~~~~~~~~~~~~~~~~~~~~~~~~~~~~~~~~~~~~~~~~~~~~~
Перша реакція екзотермічна (додає тепла, тому енергія, що виділилася в~бланкеті,
у~$M\approx1{,}1$--$1{,}3$ рази більша за~$14{,}1$~МеВ), друга, навпаки, поглинає енергію,
проте повертає нейтрон. На~один акт синтезу вигоряє один атом тритію й~народжується один
нейтрон, тож без втрат у~кращому разі вийшов би рівно один новий атом. Втрати неминучі,
тому додають \emph{розмножувач нейтронів} (берилій, свинець:
$n+{}^9\mathrm{Be}\to2n+2\,{}^4\mathrm{He}$) і~вимагають \emph{коефіцієнт відтворення
тритію} більший за~одиницю. Тепло з~бланкета забирає теплоносій (вода, гелій, рідкий
метал), далі звичайний цикл: пара, турбіна, генератор, з~термічним ККД
$\eta_{\mathrm{th}}\sim30$--$40\%$.



\subsection*{Коефіцієнт підсилення \texorpdfstring{$Q$}{Q}}



У~сталому режимі втрати плазми покриваються $\alpha$-нагрівом і~зовнішнім нагрівом:
$P_\alpha+P_{\mathrm{aux}}=P_{\mathrm{loss}}$. Означимо
% ~~~~~~~~~~~~~~~~~~~~~~~~~~~~~~~~~~~~~~~~~~~~~~~~~~~~~~~~
\begin{equation*}
	Q=\frac{P_f}{P_{\mathrm{aux}}}.
\end{equation*}
% ~~~~~~~~~~~~~~~~~~~~~~~~~~~~~~~~~~~~~~~~~~~~~~~~~~~~~~~~
Оскільки $P_\alpha=P_f/5$, частка нагріву за~рахунок $\alpha$-частинок
$P_\alpha/(P_\alpha+P_{\mathrm{aux}})=Q/(Q+5)$. При $Q=1$ вона дорівнює $1/6$
(\emph{беззбитковість} плазми: синтез дає стільки ж, скільки ми вкладаємо); при $Q=5$~---
половина; при $Q=10$~--- $2/3$; при $Q\to\infty$~--- одиниця: це і~є \emph{запалювання},
зовнішній нагрів не потрібен. Для порівняння: токамаки досягли $Q$ порядку одиниці (JET 1997: $Q\approx0{,}67$; JT-60U: еквівалентне $Q\approx1{,}25$, екстрапольоване з~D--D;
JET 2023: рекорд енергії 69~МДж, але $Q<1$; дані слід перевірити на момент видання). NIF (грудень 2022) отримав target gain $>1$,
але це ІТС і~інше означення: енергія, поглинена мішенню, а~не споживана з~мережі. Отже, «$Q$» у~різних контекстах означає різне. Промислова
електростанція потребує $Q\gg10$ (див.~нижче). Важливіше за конкретні числа~--- зв'язок $Q$ з~енергетичним
балансом.

Але $Q$~--- характеристика \emph{плазми}, а~не електростанції. Оцінимо, який $Q$ потрібен
для чистого вироблення електроенергії. Усе тепло (від нейтронів, $\alpha$-частинок
і~зовнішнього нагріву) зрештою потрапляє на стінки й~дивертор, тож
$P_{\mathrm{th}}=MP_n+P_\alpha+P_{\mathrm{aux}}$, де $P_n=0{,}8P_f$. Електричну потужність
$\eta_{\mathrm{th}}P_{\mathrm{th}}$ треба порівняти з~тією, що йде на нагрів,
$P_{\mathrm{aux}}/\eta_{\mathrm{aux}}$:
% ~~~~~~~~~~~~~~~~~~~~~~~~~~~~~~~~~~~~~~~~~~~~~~~~~~~~~~~~
\begin{equation*}
	\eta_{\mathrm{th}}\Bigl[(0{,}8M+0{,}2)\,Q+1\Bigr]\geq\frac1{\eta_{\mathrm{aux}}}.
\end{equation*}
% ~~~~~~~~~~~~~~~~~~~~~~~~~~~~~~~~~~~~~~~~~~~~~~~~~~~~~~~~~~~~~~~
При $M=1{,}2$, $\eta_{\mathrm{th}}=0{,}35$ та~$\eta_{\mathrm{aux}}=0{,}4$ (ККД систем
нагріву) це дає $Q\gtrsim5$ навіть у~ідеальному випадку. Реально треба ще живити магніти,
кріогеніку, насоси, тому проєкти електростанцій розраховують на~$Q$ порядку десятків.



\section{Що ще треба вирішити}



\begin{itemize}
	\item \emph{Вихлоп.} Гелій \enquote{попіл} накопичується й~розбавляє паливо, його
	  потрібно виводити. Для цього силові лінії на~краю відводять у~\emph{дивертор}, де
	  потік тепла сягає порядку $10$~МВт/м$^2$, тож без випромінювання домішками пластини
	  не витримали б.
	\item \emph{Матеріали.} Нейтрони $14$~МеВ зміщують атоми конструкції (десятки зміщень
	  на~атом за~рік) і~роблять її радіоактивною. Усе, що поряд із~плазмою, мусить
	  витримувати це роками.
	\item \emph{Паливний цикл.} Вигоряє лише невелика частка тритію, що подається, решту
	  потрібно уловлювати й~повертати. Тритій радіоактивний (період піврозпаду
	  $\approx12$~років).
\end{itemize}



\subsection*{Чому ця складність виправдана}



Енергія на~одиницю маси палива D--T
$17{,}6~\text{МеВ}/5\,u\approx3{,}4\times10^{14}$~Дж/кг, приблизно в~$10^{7}$ разів більше
за~вугілля. Електростанція потужністю $1$~ГВт (ел.) за~$\eta_{\mathrm{th}}\sim0{,}4$ має
$P_{\mathrm{th}}\approx2{,}5$~ГВт, тобто (за~$M=1{,}2$) $P_f\approx P_{\mathrm{th}}/(0{,}8M+0{,}2)\approx2$~ГВт, і~спалює $\approx0{,}6$~кг D\,+\,T на~добу, тоді як вугільній
потрібно близько $8$--$9$ тисяч тонн.



\section{Інерційний шлях}



Тут енергію \enquote{добувають} вибухом мікромішені, яку стискає лазерне випромінювання.
Для~оцінки успіху використовують \emph{коефіцієнт підсилення мішені} $G$~--- відношення енергії
термоядерного виходу до енергії, доставленої до мішені. Але це не те саме, що енергетичний
баланс усієї установки: лазери, нагрів, магніти, кріогеніка та інші системи також
споживають енергію. Тому слід розрізняти $Q$ плазми, підсилення мішені та інженерний
енергетичний баланс.

Мішень нагрівають \emph{прямо} (лазерні пучки б'ють по оболонці) або \emph{непрямо}
(лазер гріє стінки порожнини-хольраума, а~рентгенівське випромінювання стискає оболонку).
Стиснення має ракетний характер: абляція зовнішнього шару штовхає решту всередину.
Після стиснення центральна \emph{гаряча пляма} (hot spot) запалюється, а~навколо неї лишається
холодне щільне паливо, у~якому горіння поширюється. Аналог критерію Лоусона для ІТС~---
$\rho R\gtrsim0{,}3$~г/см$^2$: пробіг $\alpha$-частинок має бути меншим за розмір плями, а~час
розльоту $R/c_s$~--- більшим за час горіння, тобто роль добутку $n\tau$ відіграє $\rho R/c_s$.

Головна складність у~цьому підході~--- гідродинамічна нестійкість Релея~--- Тейлора
(розд.~\ref{sec:rayleigh-taylor}): під час прискореного стискання мішені важка оболонка «тисне» на легку гарячу
зону, і~нерівності поверхні зростають як $e^{\gamma t}$, $\gamma=\sqrt{Agk}$, де $A$~--- число Атвуда (розд.~\ref{sec:rayleigh-taylor}), а~роль~$g$ відіграє ефективна гравітація в~системі гальмованої оболонки, спрямована від легкого палива до важкої оболонки, руйнуючи симетрію стискання. Тому симетрія
опромінення й~гладкість мішені мають критичне значення. У~магнітному утриманні нестійкості й~перенос
розглянуто в~розд.~\ref{sec:instability} та~\ref{sec:magnetic-confinement}.

%% ========================================================
\section{Як це вимірюють: нейтронна діагностика}\label{sec:fusion-diagnostics}
%% ========================================================

Нейтрон D--T-реакції не відчуває магнітного поля й~майже не взаємодіє з~плазмою, тож вилітає з~установки та~несе інформацію про те, що відбулося всередині.
Кожен акт синтезу дає один нейтрон, тому швидкість реакції вимірюють підрахунком нейтронів поза камерою (калібровані лічильники, активаційні фольги):
% ~~~~~~~~~~~~~~~~~~~~~~~~~~~~~~~~~~~~~~~~~~~~~~~~~~~~~~~~
\begin{equation}\label{eq:neutron-yield}
	\tcbhighmath{\frac{dN_n}{dt}=\int n_Dn_T\langle\sigma v\rangle\,dV,\qquad P_f=\frac{dN_n}{dt}\cdot17{,}6~\text{МеВ}.}
\end{equation}
% ~~~~~~~~~~~~~~~~~~~~~~~~~~~~~~~~~~~~~~~~~~~~~~~~~~~~~~~~~
Отже, термоядерна потужність і~$Q=P_f/P_{\rm aux}$ визначаються прямо з~нейтронного потоку.
\emph{Оцінка.} Потужності $P_f=10$~МВт відповідає $dN_n/dt\approx3{,}5\times10^{18}$~нейтронів/с.

Енергетичний спектр дає температуру іонів: нейтрон народжується в~системі центра мас, що рухається з~тепловою швидкістю іона, тож лінія $14$~МеВ допплерівськи розширюється,
$\Delta E_{\rm FWHM}\approx82{,}5\sqrt{T_i[\text{кеВ}]}$~кеВ (при $T_i=10$~кеВ ширина $\approx260$~кеВ, порівняйте з~розд.~\ref{sec:collision-diagnostics}).

\emph{Межі застосовності.} Нейтронні лічильники потребують абсолютного калібрування з~урахуванням розсіяння в~конструкціях, а~у~дейтерієвій плазмі без тритію сигнал дають лише D--D-нейтрони ($2{,}45$~МеВ), на~кілька порядків слабші.
Допплерівську ширину на~практиці ускладнюють розсіяні нейтрони й~немаксвеллівські добавки до розподілу іонів (пучки, швидкі хвости), тому цей метод дає лише оцінку $T_i$.




%% --------------------------------------------------------


%% --------------------------------------------------------
%% --------------------------------------------------------
\begin{Summary}
\begin{itemize}
\item Найлегша реакція~--- D--T; для неї потрібна температура $T\sim10$--$20$~кеВ.
\item Пік Гамова виникає як добуток спадного максвеллівського хвоста та зростаючої ймовірності тунелювання крізь кулонівський бар'єр.
\item Критерій Лоусона зв'язує густину, температуру та~час утримання, $n\tau_E\gtrsim10^{14}$--$10^{15}~\mathrm{см}^{-3}\mathrm{с}$.
\item Випромінювальні втрати (гальмівні й~циклотронні) задають нижню межу температури запалювання.
\item Є два шляхи: магнітне утримання (великий час, мала густина) та~інерційне (мала тривалість, надвелика густина).
\item Нейтрони D--T-синтезу залишають плазму: їхній потік дає $P_f$ і~$Q$, а~допплерівська ширина лінії~--- $T_i$ (розд.~\ref{sec:fusion-diagnostics}).
\end{itemize}
\end{Summary}

\Needspace{9\baselineskip}
\section*{Задачі}
%% --------------------------------------------------------

\begin{problem}\label{prb:PlasmaFusion}
	Для дейтерій-тритієвої плазми ($n_D=n_T$, $n_e=10^{14}~\mathrm{см}^{-3}$, $k_{\mathrm
	B}T=10$~кеВ, $\langle\sigma v\rangle=1{,}1\times10^{-16}~\mathrm{см}^3/\mathrm{с}$)
	знайдіть $P_f$, $P_\alpha$, $P_{br}$ (у~Вт/$\mathrm{см}^3$) та мінімальний час
	утримання $\tau_E$ для запалювання без урахування випромінювання. Знайдіть $n\tau_E$
	при $20$~кеВ ($\langle\sigma v\rangle=4{,}3\times10^{-16}~\mathrm{см}^3/\mathrm{с}$).
	\emph{Відповідь:} $P_f\approx0{,}80$, $P_\alpha\approx0{,}16$, $P_{br}\approx0{,}017$~Вт/см$^3$; $\tau_E\approx3$~с ($n\tau_E\approx3\times10^{14}~\mathrm{см}^{-3}\mathrm c$); при 20~кеВ $n\tau_E\approx1{,}6\times10^{14}$.
\end{problem}

\begin{problem}\label{prb:AlphaHeating}
	Знайдіть швидкість і~ларморівський радіус $\alpha$-частинки з~енергією $3{,}5$~МеВ
	у~полі $5{,}3$~Тл. Покажіть, що частка нагріву плазми $\alpha$-частинками дорівнює
	$Q/(Q+5)$, і~обчисліть її для $Q=10$.
	\emph{Відповідь:} $r_{L\alpha}\approx5{,}1$~см; частка $Q/(Q+5)$, при $Q=10$ це $2/3$.
\end{problem}

\begin{problem}\label{prb:NetElectricity}
	Виведіть умову чистого вироблення електроенергії
	$\eta_{\mathrm{th}}[(0{,}8M+0{,}2)Q+1]\geq1/\eta_{\mathrm{aux}}$ і~знайдіть мінімальне
	$Q$ при $M=1{,}2$, $\eta_{\mathrm{th}}=0{,}35$, $\eta_{\mathrm{aux}}=0{,}4$. Як
	зміниться результат, якщо додатково $20\%$ електрики йде на~власні потреби (магніти,
	кріогеніка)?
	\emph{Відповідь:} $Q_{\min}\approx5{,}3$.
\end{problem}

\begin{problem}\label{prb:ICFHotSpot}
	Для інерційного синтезу задано $\rho R=0{,}3$~г/см$^2$ при $\rho=300$~г/см$^3$. Знайдіть радіус
	гарячої плями та оцініть час інерційного утримання $\tau\sim R/c_s$ при $T=10$~кеВ
	($c_s\approx\sqrt{2T/\bar m_i}$, $\bar m_i=2{,}5\,m_p$).
	\emph{Відповідь:} $R=10$~мкм, $c_s\sim9\times10^{7}$~см/с, $\tau\sim10$~пс.
\end{problem}

\begin{problem}\label{prb:DTvsDHe}
	Порівняйте реакції D--T і~D--$^3$He: енерговиділення, частку енергії в~заряджених продуктах,
	типову оптимальну температуру (рис.~\ref{fig:sigmav}). У~чому перевага й~недолік безнейтронного синтезу?
\end{problem}

\begin{problem}\label{prb:FuelConsumption}
	Яка енергія виділяється на~1~кг суміші D\,+\,T? Скільки палива на~добу споживає
	термоядерна електростанція з~електричною потужністю $1$~ГВт ($\eta=0{,}4$) і~скільки
	вугілля (теплотворність $29$~МДж/кг, $\eta=0{,}35$) потрібно звичайній?
	\emph{Відповідь:} $\approx0{,}64$~кг/добу палива проти $\approx8{,}5$~тис.~т вугілля на добу.
\end{problem}

\begin{problem}\label{prb:NeutronYield}
	Скільки D--T-нейтронів за секунду випромінює плазма з~термоядерною потужністю $P_f=500$~МВт? Яка допплерівська ширина лінії при $T_i=15$~кеВ?
	\emph{Відповідь:} $dN_n/dt=P_f/(17{,}6~\text{МеВ})\approx1{,}8\times10^{20}$~с$^{-1}$; $\Delta E\approx82{,}5\sqrt{15}\approx320$~кеВ.
\end{problem}

%% --------------------------------------------------------
%% --------------------------------------------------------

\begin{Questions}
\begin{quizlist}
\item Виведіть умову запалювання D--T-плазми. Чому існує оптимальна температура?
		  Які втрати вона не враховує?

\item У~чому відмінність між магнітним та інерційним утриманням з~погляду добутку
		  $n\tau$?

\item Як розподіляється енергія реакції D--T між $\alpha$-частинкою й~нейтроном
		  і~чому? Чим плазмовий $Q$ відрізняється від інженерного?

\item Чому реакція D--D потребує вищих температур, ніж D--T, хоча дейтерій дешевший?

\item Чим відрізняється $Q$ для NIF від $Q$ для ITER?

\end{quizlist}
\end{Questions}
